\section{Introduction}\label{sec:intro}
Under a finite evaluation budget, a black-box optimiser can improve an incumbent locally while leaving a better region unexplored. In multimodal objectives, this makes the ability to discover a promising region distinct from the ability to refine a solution within it. We examine these two aspects of performance for Bayesian optimisation (BO) and evolutionary strategies.

Global BO with expected improvement (BO-EI) selects evaluations using an acquisition function derived from a Gaussian process (GP) posterior \cite{frazier}. TuRBO-1 fits a local GP and restricts candidate generation to an adaptive trust region \cite{eriksson}. CMA-ES uses ranked objective evaluations to adapt the mean, covariance matrix and step size of a Gaussian search distribution \cite{hansen}. Independent uniform sampling provides a baseline for global coverage.

We test two hypotheses. First, after an early miss, the evaluated points of BO-EI and TuRBO-1 concentrate near the inferior minimiser. Forcing global evaluations in BO-EI or accelerating trust-region contraction in TuRBO-1 should therefore improve discovery when the available budget allows renewed exploration. Second, a CMA-ES population may contain a point below the discovery threshold while its search distribution remains too dispersed for rapid refinement. We examine how population size, parent count, and mean updates affect
settling, discovery, and precision. With the same initial mean near
the global minimiser, we expect a smaller initial step size to
improve precision attainment.  We design experiments, as seen in Sections 4 and 5, to assess these hypotheses.
